% TOP_PROBLEM: {"problemId": "problem.top500-omission-w050054-p-versus-np-conp", "canonicalTitle": "P versus NP ∩ coNP", "releaseRank": 34, "primaryDomain": "theoretical_computer_science", "primaryDomainLabel": "Theoretical computer science", "displayStatus": "open", "releaseStatus": "open"}
% !TeX program = lualatex
% Definition and sources only; this file is not an LLM solution attempt.
\documentclass[11pt]{article}
\usepackage[margin=25mm]{geometry}
\usepackage{fontspec}
\setmainfont{Latin Modern Roman}
\usepackage{amsmath,amssymb,mathtools}
\usepackage{unicode-math}
\setmathfont{Latin Modern Math}
\usepackage{xurl,hyperref}
\hypersetup{colorlinks=true,urlcolor=blue}
\setcounter{secnumdepth}{0}
\setlength{\parindent}{0pt}
\setlength{\parskip}{5pt}
\setlength{\emergencystretch}{3em}
\begin{document}
\section*{\texorpdfstring{P versus NP \ensuremath{\cap} coNP}{P versus NP ∩ coNP}}\label{34}
\subsection{Definitions and mathematical statement}
A language lies in $\mathsf{NP}\cap\mathsf{coNP}$ when both membership and nonmembership have polynomial-length certificates verifiable in deterministic polynomial time. Thus there are verifiers $V_1,V_0$ and polynomial length bounds satisfying
\[
 x\in L\iff\exists y\,V_1(x,y)=1,\qquad
 x\notin L\iff\exists z\,V_0(x,z)=1.
\]
Determine whether every such language belongs to $\mathsf P$, that is, $\mathsf P=\mathsf{NP}\cap\mathsf{coNP}$. Algorithms are uniform and unrelativized, and inputs are total languages rather than promise problems.
\par\smallskip\textit{Formulation and concept references:} \href{https://www.scottaaronson.com/democritus/lec6.html}{[S1]}.

\subsection{Short English statement}
When both answers have efficiently checkable certificates, must the answer itself be efficiently computable?
\subsection{Sources}
\begin{itemize}
\item[S1]Scott Aaronson, PHYS771 Lecture 6: P, NP, and Friends.
\url{https://www.scottaaronson.com/democritus/lec6.html}.
\textit{Formulation source.}
\item[S2]Very few problems are in NP intersect coNP but not known to be in P. What to make of that?.
\url{https://blog.computationalcomplexity.org/2024/09/very-few-problems-are-in-np-intersect.html?m=0}.
\textit{Further reference; not a proof endorsement.}
\item[S3]Random Permutations in Computational Complexity.
\url{https://arxiv.org/abs/2511.08786}.
\textit{Further reference; not a proof endorsement.}
\end{itemize}
\end{document}
